% Intended LaTeX compiler: pdflatex
\documentclass[a4paper,12pt]{article}
\usepackage[utf8]{inputenc}
\usepackage[T1]{fontenc}
\usepackage{graphicx}
\usepackage{longtable}
\usepackage{wrapfig}
\usepackage{rotating}
\usepackage[normalem]{ulem}
\usepackage{amsmath}
\usepackage{amssymb}
\usepackage{capt-of}
\usepackage{hyperref}
\usepackage{\string~/.emacs.d/common}
\author{mahmood}
\date{\today}
\title{entropy}
\hypersetup{
 pdfauthor={mahmood},
 pdftitle={entropy},
 pdfkeywords={},
 pdfsubject={},
 pdfcreator={Emacs 28.2 (Org mode 9.5.5)}, 
 pdflang={English}}
\begin{document}

\maketitle
\cite{aima2010}\\
\href{20230227001825-entropy.org}{entropy} is a measure of uncertainty of a \href{20230104221353-random_variable.org}{random variable}\\
\begin{my_example}
a random variable with only one value--a coin that always comes up heads--has no uncertainty and thus its entropy is defined as zero; thus, we gain no information by observing its value. a flip of a fair coin is equally likely to come up heads or tails, 0 or 1, this counts as “1 bit” of entropy. the roll of a fair four-sided die has 2 bits of entropy, because it takes two bits to describe one of four equally probable choices. now consider an unfair coin that comes up heads 99\% of the time. intuitively, this coin has less uncertainty than the fair coin--if we guess heads we’ll be wrong only 1\% of the time--so we would like it to have an entropy measure that is close to zero, but positive.\\
\end{my_example}
\begin{lemma}
the entropy of a random variable \(V\) with values \(v_k\), each with \href{probability.org}{probability} \(P(v_k)\), is defined as\\
\[
\text{Entropy}: \quad H(V) = \sum_{k} P(v_k)\log_2\frac{1}{P(v_k)} = -\sum_{k} P(v_k)\log_2P(v_k)
\]\\
\end{lemma}
\end{document}